
\documentclass[UTF8,12pt,a4paper]{ctexart}

\title{\textbf{Introduction}}
\date{\today}

\begin{document}
	
	\pagestyle{plain}
	
	\maketitle
	
	
	屏幕前的朋友，你好！我是肖路。这是我的第一篇 Blog。它并不涉及具体的知识或技术分享，只是想借此机会，认真介绍一下自己，也介绍一下这个刚刚建立的个人网站。
	
	我于今年六月毕业于南昌大学际銮书院综合实验班，专业是电子信息工程。今年九月，我将前往湖南大学人工智能与机器人学院攻读硕士学位，专业为人工智能。目前，我是西湖大学凌海滨老师实验室的一名访问学生。
	
	本科期间，我学习了许多电子电路、硬件以及通信相关的知识，例如模拟电子技术基础、数字电子技术基础、数字信号分析和通信原理等。不过，本科阶段最令我感兴趣、也最让我着迷的课程，却并不是本专业的课程，而是我在物理系旁听的理论力学，以及在数学系学习的概率论。
	
	我一直对较为硬核的理论知识充满兴趣。但在认真审视自己之后，我逐渐意识到，自己或许并不擅长提出理论，也不擅长进行严谨的数学分析。我真正享受的，是知识从无到有、从零开始一步步推导出来的过程。因此，我最终决定不再将纯理论研究作为自己的发展方向。
	
	过去的四年里，我经历了许多不同的尝试：从化学院的通风橱前，到制造学院的腐蚀监测实验，再到信息工程学院的电路焊接实验台；后来，我又误打误撞地来到清华园，接触到了机器人，最终走到了西湖边。这一路并不算按部就班，却十分精彩。
	
	本科阶段的学习和经历给了我许多支持，也让我有机会按照自己的兴趣去探索，而不必过分拘泥于绩点、竞赛或形式主义的活动。在际銮书院的支持和托举下，我得以顺利完成本科学业，并幸运地获得了免试攻读研究生的机会。在这里，我再次感谢南昌大学，也感谢际銮书院。
	
	这个 Blog 是我亲手用 \LaTeX{} 完成的。因为我已经很久没有认真写点什么了，而创建一个个人网站的想法，也一直存在于我的脑海之中，却始终没有真正实现。直到今天，我终于把这个想法变成了现实。
	
	建立这个网站的初衷其实很简单：我想拥有一个持续输出的机会。过去，我很喜欢在 QQ 空间、手机备忘录，或者纸质笔记本上写下自己的想法、总结以及对未来的观望。但如今，QQ 空间已经不再像从前那样活跃，我这个年龄段的大多数人也逐渐转向微信进行日常交流。手机备忘录里的文字虽然也是一种输出，却往往只是写给自己看的；至于纸质笔记本，长时间习惯了电子写作之后，我已经很难重新拿起笔，持续写下一篇长文。
	
	我始终相信，一篇好的分享，不仅能让读者有所收获，也应该成为一次自己与自己的深度对话。对于一个有些拖延、并且在许多方面的基础都还不够扎实的人来说，这个网站无疑是一个很好的开始。它既让我有机会分享自己的想法，也督促着我继续学习和进步。
	
	如果你有任何想要讨论的问题，欢迎给我发送邮件，或者添加我的微信。期待与你交流!
	
	\newpage
	\section*{}
	
	Hello to everyone reading this page. My name is Lu Xiao, and this is my first blog post. It is not intended to be a technical article or a formal piece of knowledge sharing. I simply wanted to take this opportunity to introduce myself properly and say a few words about this newly created personal website.
	
	I graduated this June from the Jiluan College Honors Program at Nanchang University, where I studied Electronic Information Engineering. This September, I will begin my master's degree in Artificial Intelligence at the College of Artificial Intelligence and Robotics, Hunan University. At the moment, I am a visiting student in Professor Haibin Ling's laboratory at Westlake University.
	
	During my undergraduate years, I studied a wide range of subjects related to electronic circuits, hardware, and communications, including Fundamentals of Analog Electronic Technology, Fundamentals of Digital Electronic Technology, Digital Signal Analysis, and Communication Principles. However, the courses that interested me most were not actually part of my major. They were Theoretical Mechanics, which I attended in the School of Physics, and Probability Theory, which I studied in the School of Mathematics.
	
	I have always been fascinated by demanding theoretical knowledge. Yet after taking an honest look at myself, I gradually realized that I was probably not particularly good at proposing theories or carrying out rigorous mathematical analysis. What I truly enjoy is following the derivation of an idea from first principles and watching knowledge take shape step by step. This realization eventually led me to decide not to pursue purely theoretical research.
	
	Over the past four years, I have tried many different things. I started at the fume hoods of the School of Chemistry, moved on to corrosion-monitoring experiments in the School of Manufacturing, and then spent time at the circuit-soldering benches in the School of Information Engineering. Later, almost by accident, I found myself standing in front of robots at Tsinghua University, and eventually arrived at West Lake. It has not been a straightforward journey, but it has certainly been an exciting one.
	
	The experiences and support I received during my undergraduate years gave me the freedom to explore according to my own interests. I did not have to be overly constrained by grades, competitions, or activities that were largely formalities. With the support of Jiluan College, I was able to complete my undergraduate studies and was fortunate enough to earn the opportunity to pursue a master's degree through recommendation. I am sincerely grateful to Nanchang University and to Jiluan College.
	
	I created this blog by hand in \LaTeX{}. I had not written anything seriously for quite some time, even though the idea of building a personal website had been in my mind for a long time. I kept putting it off, but today I finally turned that idea into something real.
	
	The original motivation for creating this website was quite simple: I wanted a place where I could regularly put my thoughts into words. I used to enjoy writing down my ideas, summaries, and thoughts about the future in QQ Space, in the notes app on my phone, or in paper notebooks. However, QQ Space is no longer as active as it once was, and most people around my age have gradually moved to WeChat for everyday communication. Notes on my phone are also a form of writing, but they are usually written only for myself. As for paper notebooks, after spending so much time writing digitally, I have found it increasingly difficult to pick up a pen and continue writing a long piece by hand.
	
	I have always believed that a good piece of writing should do more than give its readers something to take away. It should also be a deep conversation with oneself. For someone like me, who tends to procrastinate and whose foundation in many areas is still far from solid, this website is a good reason to keep moving forward. It gives me a place to share my thoughts while continuing to learn and improve.
	
	If there is anything you would like to discuss, feel free to send me an email or add me on WeChat. I look forward to hearing from you!
	
	
	
\end{document}

